\section{Methods}
\label{sec:methods}

\subsection{Participants}
\label{sec:participants}

Data were collected in two phases. In the first phase (October 2026), five participants completed the full protocol; these recordings also served to check data quality and to fix the remaining protocol parameters. In the second phase (November and December 2026), about ten further participants were recruited with the aim of covering a range of expertise with the two movements, from stick-naive adults to practitioners with regular practice \confirm{final N and recruitment channels}. The target sample was therefore $N \approx 15$.

Inclusion criteria were age between 18 and 45 years, no known cardiovascular, respiratory or neurological condition, and the ability to perform both movements with each hand for two minutes. Stick-naive participants received a short standardised instruction and practice period before recording (Section~\ref{sec:paradigm}). Exclusion criteria were arrhythmia, an implanted cardiac device, pregnancy, and medication known to affect heart rate or HRV, such as beta-blockers. Participants were asked to avoid caffeine for two hours and strenuous exercise for 24 hours before the session, and to avoid a heavy meal in the preceding two hours \citep{laborde2017}. Compliance with these requests, hours of sleep the previous night and handedness were recorded on arrival.

Expertise was assessed with a short questionnaire covering years of practice with the stick or comparable flow tools (staff, poi, juggling, contact props), typical practice hours per week, and a self-rating of proficiency for each movement on a scale from 1 to 10 \confirm{expertise items and grouping rule}. Participants were assigned to a novice or an experienced group by a rule fixed before recruitment of the second phase, and the continuous measures were retained for a secondary analysis (Section~\ref{sec:stats}).

All participants gave written informed consent. \todo{ethics: name the approving committee or state that the study falls under the unit's existing approval; add reference number}

With a paired design, $N = 15$ provides 80\% power to detect a within-participant effect of $d \approx 0.78$ at $\alpha = .05$ (two-sided). Effects smaller than this, and all between-group comparisons of expertise, are therefore reported as estimates with confidence intervals rather than as decisive tests.

\subsection{The Stick and the Movements}
\label{sec:stick}

The stick is about 50~cm long with a weight at each end \confirm{exact length, mass and material}. The weighting moves the centre of mass outwards and gives the stick rotational inertia, so that once set in motion it tends to continue in a figure-eight path around the hand. The design was developed within a movement practice that treats structured variability, neither rigid repetition nor random movement, as the basis of skilled coordination \confirm{whether to cite the EightOS / seks.design material here}.

Two families of movement were recorded. In an \emph{isolation}, one point of the stick, typically its centre or one end, is held fixed in space while the hand and the rest of the stick move around it. Isolations are slow, demand continuous visual monitoring and fine force control, and break down visibly when control lapses. In a \emph{spin}, the stick rotates continuously around the hand or fingers in a figure-eight. Spins are faster, rely more on timing and momentum than on visual correction, and produce a regular rhythm once established. Each movement was performed with the left and with the right hand. Participants were instructed to keep the movement going for the whole trial at a pace they could sustain, to restart immediately after a drop, and to breathe naturally. Drops were counted by the experimenter.

\subsection{Experimental Paradigm}
\label{sec:paradigm}

The study used a within-participant design with one factor, \emph{condition}, at three levels (seated rest, isolation, spin), crossed with the between-participant factor \emph{expertise} (novice, experienced). Hand (left, right) was recorded within each movement condition and treated as a secondary factor. Seated rest without the stick served as the control condition; each participant was their own control.

Each session followed the sequence in Table~\ref{tab:session}. After consent, questionnaires and fitting of the chest strap, participants sat quietly for a pre-task rest period. The four movement trials (isolation left, isolation right, spin left, spin right) followed in an order counterbalanced across participants \confirm{counterbalancing scheme: Latin square vs randomised}. Each movement trial was followed by seated recovery, so that every trial started from a comparable resting state rather than from the arousal left by the previous one. A post-task rest period closed the session.

\begin{table}[ht]
\centering
\caption{Session timeline. Durations in square brackets are still to be fixed from the first-phase data.}
\label{tab:session}
\small
\begin{tabularx}{\textwidth}{l l X}
\toprule
Step & Duration & Content \\
\midrule
1 & 10 min & Consent, questionnaires (demographics, handedness, expertise, caffeine, sleep), strap fitting, signal check \\
2 & 5 min & Familiarisation: practice of both movements with each hand (novices); warm-up (experienced) \\
3 & [2--5 min] & Pre-task rest, seated, eyes open, no stick \\
4 & 4 $\times$ 2 min & Movement trials (isolation and spin, each hand), counterbalanced order \\
  & [2 min] each & Seated recovery after each movement trial \\
5 & [2--5 min] & Post-task rest, seated \\
6 & 5 min & Post-session ratings and debrief \\
\midrule
 & $\approx$ 45 min & Total \\
\bottomrule
\end{tabularx}
\end{table}

The trial duration of two minutes was chosen as a compromise between two constraints. \al{} is conventionally estimated from a two-minute window (Section~\ref{sec:bg-dfa}), so shorter trials would not yield a stable estimate. Longer trials, in turn, exceed the duration for which isolations in particular can be sustained without fatigue of the forearm and repeated drops, which would turn the trial into a series of restarts. \confirm{2 min vs 3 min: if trials are extended to 3 min, \al{} is taken from the final 2 min so that the onset transient is excluded}

Rest periods were recorded with the same capture system and the same durations as the analysis windows, so that rest and movement estimates were based on comparable amounts of data. The pre-task rest is the baseline for H1; the post-task rest serves the carry-over analysis.

\subsection{Technical Setup}
\label{sec:setup}

Figure~\ref{fig:setup} gives an overview of the recording chain. RR intervals were recorded with a Polar H10 chest strap (Polar Electro, Kempele, Finland) worn at the level of the lower sternum with moistened electrodes. The H10 transmits RR intervals at a resolution of 1/1024~s through the standard Bluetooth Low Energy Heart Rate Measurement characteristic, in which a single notification can carry several RR values.

\begin{figure}[ht]
\centering
\fbox{\parbox{0.9\textwidth}{\centering\vspace{1.5cm}
Polar H10 \quad$\rightarrow$\quad SomaSync (iPhone, BLE) \quad$\rightarrow$\quad WebSocket (local network) \quad$\rightarrow$\quad capture server (laptop) \quad$\rightarrow$\quad JSON per trial \quad$\rightarrow$\quad offline pipeline
\vspace{1.5cm}}}
\caption{Recording chain. \todo{replace with a drawn diagram and a photo of the setup}}
\label{fig:setup}
\end{figure}

A first version of the setup connected the H10 directly to a browser page through the Web Bluetooth interface. On the laptop used for recording, this connection repeatedly dropped within seconds of being established, so the final setup routed the signal through a phone. The H10 was paired with the SomaSync application on an iPhone \confirm{app version}, which maintains the Bluetooth connection natively and streams each incoming RR interval, together with its own real-time HRV metrics, as JSON messages over a WebSocket connection on the local network.

A custom capture server (Node.js) running on the laptop received this stream. Its web interface, opened in a browser on the same laptop, allowed the experimenter to enter the participant identifier, the condition label and the trial duration before starting each trial. The server stamped each received RR interval with the time of arrival, kept a running count of samples and stopped the trial automatically when the set duration had elapsed, so that every trial had the same length without manual timing. Each trial was saved as a JSON file named from the condition, participant and start time, containing the list of RR intervals with timestamps, the trial metadata and the target duration, together with a raw log of every message received. Keeping the raw log means that any value reported by the streaming application can be recomputed or checked later. The code for the capture server and the analysis pipeline is available in the project repository \todo{repository link or archive DOI}.

\subsection{Real-time Metrics}
\label{sec:realtime}

In addition to the raw RR intervals, SomaSync reports its own running estimates of heart rate, RMSSD, SDNN, pNN50 and \al, which were stored with each trial. The settings behind these estimates (window length, detrending, artifact handling) are not documented in enough detail to be reproduced \confirm{check with SomaSync documentation or developer}. They were therefore not used for the primary analyses. Instead, they were compared with the offline estimates computed from the same RR intervals (Section~\ref{sec:stats}), which tests whether a consumer application can stand in for an offline pipeline in future, larger studies.

\subsection{Questionnaires}
\label{sec:questionnaires}

Before the session, participants completed questions on demographics, handedness, expertise (Section~\ref{sec:participants}), caffeine intake, sleep and recent exercise. After each movement trial, participants rated perceived difficulty and perceived exertion \confirm{scales: e.g. Borg CR10 for exertion}, and after the session they completed a short flow questionnaire for each movement \confirm{instrument, e.g. Flow Short Scale}. These ratings are used descriptively and in exploratory analyses only.
